\section{Cenni Teorici}
\label{sec:Teoria}
\thispagestyle{plain}
In questa sezione vengono riportati alcuni cenni teorici sui concetti fondamentali dell'analisi dei segnali, che sono stati utilizzati durante lo sviluppo del progetto. Verranno trattati argomenti come la rappresentazione dei segnali, la trasformata di Fourier, il campionamento e la ricostruzione dei segnali, nonché alcune tecniche di filtraggio e analisi spettrale. Questi concetti forniscono le basi teoriche necessarie per comprendere le operazioni svolte sui segnali durante l'esperimento e per interpretare i risultati ottenuti.